\documentclass[11pt]{article}

\usepackage[margin=1.15in]{geometry}
\usepackage{amsmath,amssymb,amsthm}
\usepackage[hidelinks]{hyperref}

\DeclareMathOperator*{\argmax}{arg\,max}
\DeclareMathOperator*{\argmin}{arg\,min}
\newcommand{\E}{\mathbb{E}}
\newcommand{\R}{\mathbb{R}}
\newcommand{\N}{\mathbb{N}}
\newcommand{\1}{\mathbf{1}}
\newcommand{\DeltaK}{\Delta^k}
\newcommand{\OmegaN}{\Omega}

\newtheorem{theorem}{Theorem}
\newtheorem{proposition}{Proposition}
\newtheorem{lemma}{Lemma}
\newtheorem{corollary}{Corollary}
\newtheorem{remark}{Remark}
\theoremstyle{definition}
\newtheorem{definition}{Definition}
\theoremstyle{plain}

\title{The Informational Efficiency of Frequency-Report Scoring Rules}
\author{Jesper Armouti-Hansen}
\date{\today}

\begin{document}

\maketitle

\begin{abstract}
In a multinomial frequency-report environment, the subject has latent beliefs \(p\), but the researcher observes only an incentivized count report \(r\).
This paper studies what can be inferred from that report.
Each scoring rule induces an optimal-report correspondence and therefore an inverse belief region: the set of beliefs under which the observed report would have been optimal.
The paper evaluates the informational efficiency of such mechanisms through coordinate widths, functional widths, implementation cost, and robustness to risk aversion.
The main analytical result is that quadratic-distance frequency scoring gives a transparent projection rule, a linear belief region, closed-form coordinate bounds, and linear-program bounds for means.
Discrete-metric frequency scoring is used as the known fixed-prize rule with analytic bounds, Manhattan distance is included as a structured comparison rule, and Hamming and Chebyshev distance are discussed as secondary computable rules.
The contribution is a finite-sample design guide for turning frequency reports into belief and mean bounds.
\end{abstract}

\section{Introduction}

Belief elicitation in experiments often asks subjects to report probabilities directly.
Frequency reports ask instead for counts: out of \(n\) realizations, how many times will each outcome occur?
This format is natural when the object of interest is a distribution of choices, signals, or outcomes in a finite group.
It may also be easier for subjects than reporting abstract probabilities under a standard scoring rule.
Many probability scoring rules, including the quadratic scoring rule, are well understood theoretically \cite{Savage1971,GneitingRaftery2007,Selten1998}, but they require subjects to understand how probability reports map into payments.
That mapping can be cognitively demanding, especially when subjects reason more naturally in frequencies than in abstract probabilities \cite{Hogarth1975,SchlagEtAl2013}.
Direct monetary scoring rules also raise the familiar risk-aversion problem \cite{ArmantierTreich2012}.

The practical problem is that the researcher observes a count report \(r\), not the subject's latent belief vector \(p\).
The report is an incentivized message generated by a scoring rule.
Thus the inferential question is inverse: given the rule and the observed report, which beliefs are consistent with optimal reporting?
The paper follows the chain
\[
\text{frequency-report scoring rule}
\rightarrow
R_S(p)
\rightarrow
P_S(r)
\rightarrow
\text{belief and functional bounds}.
\]
Here \(R_S(p)\) is the set of reports that are optimal for a subject with beliefs \(p\), and
\[
P_S(r)=\{p\in\Delta^k:r\in R_S(p)\}
\]
is the inverse belief region induced by the observed report.
From this set, the researcher can derive bounds on individual probabilities and on functionals such as expected values.

This paper uses the term informational efficiency in an operational sense.
A rule is more informative for a given design objective if the observed report leads to narrower coordinate bounds or narrower mean and functional bounds.
These inferential gains must be weighed against implementation costs: exact-match payment probabilities, cognitive simplicity, computational burden, and robustness to risk aversion.

The paper's main analytical contribution is for quadratic-distance frequency scoring.
Under this rule, a subject reports the feasible count vector closest to the expected count vector \(np\).
This gives a linear belief region and closed-form bounds for each \(p_i\).
Mean bounds and other linear functional bounds are then simple linear programs over the same region.

The paper compares this rule to several other count-loss mechanisms.
Schlag--Tremewan frequency guessing is the known discrete-metric exact-match mechanism: it elicits multinomial modes, gives closed-form bounds, and is robust to risk aversion under a fixed-prize implementation \cite{SchlagTremewanSimple}.
Manhattan-distance scoring is included as a median-oriented comparison with an explicit threshold representation.
Hamming and Chebyshev rules are discussed as secondary computable rules because their inverse regions are exactly specified by finite expected-loss inequalities, but they are not the focus of the headline design exercise.

The contribution is therefore not a new exact-match mechanism and not a general theory of scoring rules.
It is a finite-sample method for interpreting frequency reports as belief and mean bounds, together with a design comparison of simple mechanisms that differ in transparency, robustness, and inferential precision.

\section{Setup}

The environment is a finite belief-elicitation problem.
There are \(k\) possible outcomes, and the subject has latent beliefs
\[
p=(p_1,\ldots,p_k)\in\DeltaK,
\quad
p_i\ge 0,
\quad
\sum_{i=1}^k p_i=1.
\]
The elicitation task asks about \(n\in\N\) independent repetitions.
If the subject's beliefs are \(p\), the realized count vector is
\[
\omega=(\omega_1,\ldots,\omega_k)\sim \mathrm{Mult}(n,p).
\]
The set of feasible count vectors is
\[
\OmegaN=\left\{\omega\in\N_0^k:\sum_{i=1}^k \omega_i=n\right\},
\]
and the multinomial probability mass function is
\[
\Pr_p(\omega)
=
\frac{n!}{\omega_1!\cdots \omega_k!}
\prod_{i=1}^k p_i^{\omega_i}.
\]
The subject reports another feasible count vector
\[
r=(r_1,\ldots,r_k)\in\OmegaN.
\]

\begin{definition}[Frequency-report scoring rule]
A frequency-report scoring rule is a function
\[
S:\OmegaN\times\OmegaN\to\R,
\]
where \(S(r,\omega)\) is the score assigned to report \(r\) when the realized count vector is \(\omega\).
The paper studies rules that reward reports close to realized counts.
A convenient form is
\[
S_D(r,\omega)=a-bD(r,\omega),
\quad b>0,
\]
where \(D:\OmegaN\times\OmegaN\to\R_+\) is a count-error index.
The function \(D\) need not be a metric: squared distance is a divergence-style count loss, not a metric, because it fails the triangle inequality.
\end{definition}

The constants \(a\) and \(b\) do not matter for expected-score rankings as long as \(b>0\).
They matter for implementation.
If scores are paid directly as money and nonnegative payments are desired, one can choose
\[
a\ge b\max_{r,\omega\in\OmegaN}D(r,\omega).
\]
How the score is implemented as a payment affects robustness to the subject's risk attitude; the paper's discussion of risk aversion and implementation takes up this question separately.

\begin{definition}[Optimal reports]
Under risk neutrality, the optimal-report correspondence induced by a scoring rule \(S\) is
\[
R_S(p)
=
\argmax_{r\in\OmegaN}\E_p[S(r,\omega)].
\]
Because \(\OmegaN\) is finite, \(R_S(p)\) is nonempty.
\end{definition}

Risk neutrality is the maintained assumption throughout the body of the paper; the discussion of risk aversion and implementation separately shows how the analysis extends to risk-averse expected-utility subjects.

This is the forward incentive object.
It answers the question: if the subject's latent beliefs are \(p\), which reports are optimal?
The econometric object is the inverse question.
After observing \(r\), the researcher asks which beliefs could have generated that report.

\begin{definition}[Identified set]
For an observed report \(r\in\OmegaN\), the identified set induced by \(S\) is
\[
P_S(r)
=
\{p\in\DeltaK:r\in R_S(p)\}.
\]
Ties are included: if \(r\) is one of several optimal reports under \(p\), then \(p\in P_S(r)\).
\end{definition}

We call \(P_S(r)\) the \emph{mechanism-induced identified set}: the identifying restriction is the optimal-report condition \(r\in R_S(p)\) imposed by the scoring-rule mechanism, not a sampling distribution over repeated observations. By construction \(P_S(r)\) is exactly the set of beliefs consistent with that restriction.

The most direct summaries of \(P_S(r)\) are coordinate-wise identification bounds:
\[
\underline p_{S,i}(r)
=
\inf_{p\in P_S(r)}p_i,
\quad
\overline p_{S,i}(r)
=
\sup_{p\in P_S(r)}p_i,
\]
with interval
\[
P_{S,i}(r)
=
[\underline p_{S,i}(r),\overline p_{S,i}(r)].
\]
When the scoring rule is clear from context, the subscript \(S\) is omitted.

Many experiments require more than coordinate-by-coordinate probability bounds.
If outcomes have numerical values \(x_1,\ldots,x_k\), the latent mean implied by beliefs \(p\) is
\[
\mu(p;x)=\sum_{i=1}^k p_i x_i.
\]
This could represent an expected outcome, an expected treatment share, a predicted group average, or another linear index of the belief distribution.
The sharp mean bounds after observing \(r\) are
\[
\underline\mu_S(r;x)
=
\inf_{p\in P_S(r)}\sum_{i=1}^k p_i x_i,
\quad
\overline\mu_S(r;x)
=
\sup_{p\in P_S(r)}\sum_{i=1}^k p_i x_i.
\]
The corresponding mean interval is
\[
[\underline\mu_S(r;x),\overline\mu_S(r;x)].
\]
These bounds generally require optimizing over the full joint region \(P_S(r)\).
Coordinate intervals alone do not generally give sharp bounds for a mean, because they ignore dependence across coordinates imposed by the simplex and the scoring rule.

The design comparison later summarizes the informativeness of a rule using coordinate widths,
\[
\overline p_{S,i}(r)-\underline p_{S,i}(r),
\]
and functional widths,
\[
\overline\mu_S(r;x)-\underline\mu_S(r;x).
\]
Exact-match payment probability is reported only for the discrete-metric rule, as an implementation diagnostic of that fixed-prize mechanism; it is not a cross-rule comparison metric.

The next section applies this setup to the three headline rules: quadratic-distance, discrete-metric, and Manhattan.

\section{Frequency-Report Scoring Rules}

This section applies the setup to five count-loss rules.
The first two rules have simple analytic identification bounds and are the primary theoretical objects in the paper.
The remaining rules have exact inverse-region representations that are useful for computational comparison, but their sharp bounds are obtained from those representations rather than from comparably simple main-text formulas.

\subsection{Primary Rules with Analytic Bounds}

\subsubsection{Quadratic-Distance Frequency Scoring}

The main analytical rule is quadratic-distance frequency scoring.
It pays less when the reported count vector is far from the realized count vector in squared Euclidean distance:
\[
S_Q(r,\omega)=a-b\|r-\omega\|_2^2,
\quad b>0.
\]
This rule is connected in spirit to quadratic scoring rules, but it is not a standard probability score.
It is a frequency-report count-loss rule: subjects report counts, and the penalty is the squared difference between reported and realized counts.
Write \(P_Q\) for the belief region induced by this rule.

Given beliefs \(p\), the expected count vector is \(np\).
The key forward fact is that quadratic-distance scoring asks the subject to report the feasible integer count vector closest to \(np\).
The inverse question is therefore geometric: for which beliefs is the observed report \(r\) the closest feasible count vector to \(np\)?

\begin{proposition}[Quadratic distance: reports, regions, and bounds]
For every \(p\in\DeltaK\),
\[
R_{S_Q}(p)
=
\argmin_{s\in\OmegaN}\|s-np\|_2^2.
\]
For an observed report \(r\in\OmegaN\), the belief identification region is
\[
P_Q(r)
=
\left\{
p\in\DeltaK:
n(p_i-p_j)\le r_i-r_j+1
\quad
\forall i,j\text{ with }r_j>0
\right\}.
\]
Let
\[
m(r)=|\{j:r_j>0\}|
\]
be the number of categories with positive reported counts.
The sharp coordinate intervals are
\[
P_{Q,i}(r)=
\begin{cases}
\displaystyle
\left[0,\frac{m(r)}{n(m(r)+1)}\right],
& r_i=0,\\[1.2em]
\displaystyle
\left[
\frac{r_i-1}{n}+\frac{1}{nk},
\frac{r_i+1}{n}-\frac{1}{n\,m(r)}
\right],
& r_i>0.
\end{cases}
\]
Projection ties are included.
\end{proposition}

The inequalities compare the observed report \(r\) with reports obtained by moving one count from a positive reported category \(j\) to another category \(i\).
If all such one-count moves weakly increase squared distance from \(np\), then \(r\) is a closest feasible report.
Thus \(P_Q(r)\) is a polytope in the belief simplex.
The proof is in Appendix \ref{app:squared-proof}.

\paragraph{Probability and mean bounds.}
The coordinate bounds are closed form.
For numerical outcomes \(x_1,\ldots,x_k\), sharp mean bounds are obtained by optimizing over the same polytope:
\[
\min_{p\in P_Q(r)}\sum_{i=1}^k p_i x_i
\quad\text{and}\quad
\max_{p\in P_Q(r)}\sum_{i=1}^k p_i x_i.
\]
These are linear programs.
For coordinate probabilities, the displayed intervals give the corresponding linear-program solutions explicitly.
This is the main practical advantage of quadratic-distance frequency scoring: the same transparent region gives both probability bounds and bounds on expected values.

\paragraph{Practical interpretation.}
Quadratic distance is mean-oriented: it links a report \(r\) to beliefs for which \(r\) is the integer projection of \(np\).
Compared with exact matching, the payment varies smoothly with distance rather than only with exact equality.
The main implementation cost is that direct monetary quadratic-distance payments generally require risk neutrality unless the score is implemented probabilistically, as discussed below.

\subsubsection{Discrete-Metric Frequency Scoring}

The second primary rule is the frequency-guessing mechanism of Schlag and Tremewan \cite{SchlagTremewanSimple}.
It is the count-vector discrete metric:
\[
D_0(r,\omega)=\1\{r\ne\omega\}.
\]
Equivalently, a fixed prize is paid if and only if the realized count vector is exactly equal to the report:
\[
q_0(r,\omega)=\1\{r=\omega\}.
\]
Under risk neutrality, this is equivalent to the affine score
\[
S_0(r,\omega)=a+b\1\{r=\omega\},
\quad b>0.
\]
The fixed-prize formulation is useful for implementation and risk aversion; the affine-score formulation is convenient for comparison with the distance rules.
Write \(R_0\) and \(P_0\) for the report correspondence and belief region induced by this rule.

The mechanism is simple: a subject with beliefs \(p\) maximizes the probability of an exact match,
\[
\Pr_p(\omega=r)
=
\frac{n!}{r_1!\cdots r_k!}\prod_{i=1}^k p_i^{r_i}.
\]
Thus the forward object is the multinomial mode set.
The inverse object is the set of beliefs for which the observed report is one of those modes.

\begin{proposition}[Discrete metric: reports, regions, and bounds]
Let \(r\in\OmegaN\), and allow ties in the multinomial mode.
Then \(r\) is optimal under discrete-metric frequency scoring if and only if \(r\) is a multinomial mode:
\[
R_0(p)
=
\argmax_{s\in\OmegaN}\Pr_p(\omega=s).
\]
The belief identification region is
\[
P_0(r)
=
\left\{
p\in\DeltaK:
r_jp_i\le (r_i+1)p_j
\quad
\forall i,j\text{ with }r_j>0
\right\}.
\]
For each coordinate, the sharp identification interval is
\[
P_{0,i}(r)
=
\left[
\frac{r_i}{n+k-1},
\frac{r_i+1}{n+1}
\right].
\]
\end{proposition}

The inequalities say that no one-count transfer from a reported category \(j\) to another category \(i\) can increase the multinomial probability of the report.
Weak inequalities include ties.
The formula also handles boundaries: if \(r_i=0\), the lower bound is \(0\); if \(r_i=n\), the upper bound is \(1\).
These are the frequency-guessing bounds in Schlag and Tremewan \cite[Proposition 1]{SchlagTremewanSimple}; they are restated here to make the informational-efficiency comparison self-contained.
The proof is in Appendix \ref{app:exact-match-proof}.

\paragraph{Probability and mean bounds.}
The discrete-metric rule gives closed-form coordinate intervals.
For linear functionals such as means, sharp bounds use the full joint region, not just the coordinate intervals:
\[
\min_{p\in P_0(r)}\sum_{i=1}^k p_i x_i
\quad\text{and}\quad
\max_{p\in P_0(r)}\sum_{i=1}^k p_i x_i.
\]
These are linear programs because \(P_0(r)\) is defined by linear inequalities.
Thus discrete-metric frequency scoring is analytically convenient for individual probability bounds, while mean inference still requires using the joint mode region.

\paragraph{Practical interpretation.}
Discrete-metric frequency scoring is transparent and robust to risk aversion because it is a fixed-prize event.
Its implementation cost is that the probability of an exact match can be small when \(n\) is large or the outcome space is rich.
It is a prominent known rule in the comparison, not a new mechanism introduced here.

\subsection{Additional Rules for Computational Comparison}

The next rules broaden the design exercise.
For each rule, the inverse region is fully specified by finite expected-loss inequalities and can therefore be used in the simulation exercise.
What distinguishes these rules from the primary rules is not lack of identification, but the absence of a simple main-text formula for sharp coordinate and mean bounds comparable to the quadratic-distance and discrete-metric cases.
For these rules, the computational representation is the relevant object for the design exercise.

\subsubsection{Manhattan-Distance Frequency Scoring}

The Manhattan-distance rule is
\[
S_M(r,\omega)=a-b\sum_{i=1}^k |r_i-\omega_i|,
\quad b>0.
\]
It rewards total absolute count accuracy rather than squared count accuracy.
Its inferential interpretation is median-oriented rather than mean-oriented.

\paragraph{Optimal reports and belief region.}
For category \(i\), the realized count has binomial marginal distribution
\[
W_i\sim\mathrm{Bin}(n,p_i).
\]
Let
\[
\ell_i(t;p_i)=\E|t-W_i|,
\qquad
F_i(t;n,p_i)=\Pr(W_i\le t).
\]
Then optimal reports solve
\[
R_M(p)
=
\argmin_{r\in\OmegaN}
\sum_{i=1}^k \ell_i(r_i;p_i).
\]
The simplex constraint \(\sum_i r_i=n\) prevents the researcher from treating each coordinate independently.
The relevant optimality condition is that no one-count transfer lowers expected Manhattan loss.
Equivalently,
\[
P_M(r)
=
\left\{
p\in\DeltaK:
F_i(r_i;n,p_i)\ge F_j(r_j-1;n,p_j)
\quad
\forall i,j\text{ with }r_i<n,\ r_j>0
\right\}.
\]
Weak inequalities include ties; there is no receiver constraint when \(r_i=n\) and no sender constraint when \(r_j=0\).

\paragraph{Identification bounds.}
For \(k=2\), write \(r=(t,n-t)\).
Then Manhattan loss is \(2|t-W|\), with \(W\sim\mathrm{Bin}(n,p_1)\), so the report identifies the binomial-median interval
\[
P_{M,1}(r)
=
\{p_1\in[0,1]:F(t-1;n,p_1)\le 1/2\le F(t;n,p_1)\}.
\]
This binary case is clean.

For \(k>2\), the CDF inequalities remain explicit.
The useful computational representation is a one-dimensional threshold search.
Specifically, \(p\in P_M(r)\) if and only if there exists \(c\in[0,1]\) such that
\[
F_i(r_i-1;n,p_i)\le c\le F_i(r_i;n,p_i)
\quad\forall i,
\]
using \(F_i(-1;n,p_i)=0\) and \(F_i(n;n,p_i)=1\).
For a fixed \(c\), these inequalities give coordinate boxes through inverse binomial CDFs; intersecting those boxes with \(\sum_i p_i=1\) gives the feasible beliefs at that threshold.
The sharp coordinate and mean bounds are computed by optimizing over feasible thresholds.

\paragraph{Practical interpretation.}
Manhattan distance is useful when a median-type count prediction is substantively appropriate.
It is also useful as a contrast: changing the distance criterion changes the inferred belief region.
For mean-oriented belief elicitation, however, the multi-category Manhattan bounds are less transparent than the quadratic-distance bounds.

\subsubsection{Hamming-Distance Frequency Scoring}

The Hamming-distance rule penalizes the number of categories whose reported count differs from the realized count:
\[
S_H(r,\omega)=a-b\sum_{i=1}^k\1\{r_i\ne\omega_i\},
\quad b>0.
\]
Its expected loss separates into binomial marginal exact-coordinate probabilities:
\[
\E_p\left[\sum_i \1\{r_i\ne\omega_i\}\right]
=
\sum_i \{1-\Pr(\mathrm{Bin}(n,p_i)=r_i)\}.
\]
Thus the forward objective is computationally simple, even though the reported counts still have to sum to \(n\).

For an observed report \(r\), define
\[
L_H(r;p)
=
\sum_i \{1-\Pr(\mathrm{Bin}(n,p_i)=r_i)\}.
\]
The inverse region is exactly
\[
P_H(r)
=
\{p\in\DeltaK:L_H(r;p)\le L_H(s;p)\ \forall s\in\OmegaN\}.
\]
This finite set of expected-loss inequalities fully specifies identification.
Unlike the primary rules, however, the region is not a simple transfer polytope or a closed coordinate interval formula.
The bounds are therefore computed from the exact inverse-region inequalities in the design exercise.

For \(k=2\), Hamming distance coincides with discrete-metric frequency scoring.
Matching one coordinate is equivalent to matching the whole binary count vector, so the report incentives are the same.

\subsubsection{Chebyshev-Distance Frequency Scoring}

The Chebyshev-distance rule penalizes the largest coordinate-level count error:
\[
S_\infty(r,\omega)=a-b\max_i |r_i-\omega_i|,
\quad b>0.
\]
For \(k=2\), with \(r=(t,n-t)\) and \(\omega=(W,n-W)\),
\[
\|r-\omega\|_\infty=|t-W|.
\]
Manhattan loss is \(2|t-W|\) in the same binary environment.
Thus Chebyshev and Manhattan distance induce the same optimal reports when \(k=2\).

For \(k>2\), the maximum couples coordinates inside each realization.
The expected-loss comparison is therefore nonseparable.
The inverse region is exactly
\[
P_\infty(r)
=
\left\{
p\in\DeltaK:
\E_p\|r-\omega\|_\infty
\le
\E_p\|s-\omega\|_\infty
\quad
\forall s\in\OmegaN
\right\}.
\]
This region is finite and computationally checkable by enumerating feasible reports and evaluating expected losses.
As with Hamming distance, the bounds used in the design exercise are computed from the exact inverse-region inequalities rather than from a simple closed main-text formula.

\section{Informational Efficiency Exercise}

The analytical regions are useful because they tell an experimenter what can be inferred after observing a report.
This section reports a design exercise generated by the reproducible script \texttt{scripts/design\_efficiency.py}.
The exercise draws latent beliefs
\[
p\sim\operatorname{Dirichlet}(\alpha,\ldots,\alpha)
\]
and, for each rule, computes an optimal report \(r_S(p)\) and the bounds induced by \(P_S(r_S(p))\).
The final run uses \(5{,}000\) belief draws for each cell in
\[
n\in\{5,10,20,50\},\qquad
k\in\{2,3,5,10\},\qquad
\alpha\in\{0.1,0.3,1,3,10\}.
\]
Manhattan bounds are threshold-computed with tolerance \(10^{-4}\).

The comparison is conditional on the theoretical characterizations above.
It is not evidence of incentive compatibility, and it is not an optimality theorem.
It is a diagnostic for the practical question: which rule gives a researcher tighter inference for the object of interest?

\paragraph{Aggregate comparison.}
Tables \ref{tab:design-widths} and \ref{tab:design-wins} report averages over the full design grid.
Lower widths indicate more precise inference.
Win shares report the average share of belief draws for which a rule gives the narrowest interval for the stated target.

\begin{table}[h]
\centering
\small
\begin{tabular}{lrrrr}
\hline
Rule & avg coord & max coord & linear mean & skewed mean\\
\hline
Discrete metric & 0.1025 & 0.1478 & 0.1142 & 0.1025\\
Quadratic distance & 0.1042 & 0.1185 & 0.1112 & 0.1042\\
Manhattan distance & 0.1021 & 0.1253 & 0.1101 & 0.1021\\
\hline
\end{tabular}
\caption{Average widths in the latent-belief design exercise. Columns average over all \((n,k,\alpha)\) cells. Manhattan entries are threshold-computed.}
\label{tab:design-widths}
\end{table}

\begin{table}[h]
\centering
\small
\begin{tabular}{lrrr}
\hline
Rule & avg-coordinate wins & max-coordinate wins & linear-mean wins\\
\hline
Discrete metric & 0.6204 & 0.3431 & 0.5266\\
Quadratic distance & 0.3189 & 0.6070 & 0.3776\\
Manhattan distance & 0.0607 & 0.0499 & 0.0958\\
\hline
\end{tabular}
\caption{Average win shares in the latent-belief design exercise. A win is the smallest interval width among the three rules for a draw and target. The linear mean uses \(x_i=(i-1)/(k-1)\).}
\label{tab:design-wins}
\end{table}

The aggregate results do not support a universal ranking.
Discrete-metric scoring gives the narrowest average coordinate interval most often.
Quadratic-distance scoring is strongest for worst-coordinate uncertainty, with a win share of about \(0.607\) for maximum coordinate width.
It is also competitive for mean inference and performs especially well in sparse-belief environments.
Manhattan distance is a useful structured comparison rule, but this run does not support presenting it as the dominant design choice.

\paragraph{Where quadratic distance helps.}
The cell-level results show that quadratic distance is particularly useful when beliefs are sparse or concentrated.
For example, with \(\alpha=0.1\), quadratic distance wins the linear-mean criterion in more than \(0.8\) of draws for many cells, including \(n=50,k=10\), \(n=50,k=5\), and \(n=20,k=5\).
It also wins the worst-coordinate criterion almost always in several larger-\(k\) cells.
This is the setting in which the projection interpretation is most valuable: the rule translates reports near sparse or extreme count vectors into tight bounds on high-probability categories and mean-oriented functionals.

\paragraph{Payment probability.}
The discrete-metric rule has a fixed-prize implementation, which is attractive under risk aversion.
Its implementation cost is that the exact-match payment probability can be small.
In the final run, the average exact-match probability is high for very sparse binary beliefs, but it falls sharply as \(n\), \(k\), and belief balance increase; for instance, it is essentially zero in the \(n=50,k=10,\alpha\in\{1,3,10\}\) cells.
Thus narrow bounds under discrete-metric scoring should be weighed against the chance that the subject rarely wins the fixed prize.

\paragraph{Fixed-report illustrations.}
Table \ref{tab:selected-design} holds the observed report fixed and compares the induced bounds.
This answers a different question from the latent-belief simulation: conditional on seeing the same report, how does the scoring rule change the inferred belief set?

\begin{table}[h]
\centering
\small
\begin{tabular}{llrrr}
\hline
Report and rule & payment prob. & avg coord width & linear mean interval & method\\
\hline
\multicolumn{5}{l}{\(n=10,\ k=3,\ r=(5,3,2)\)}\\
Discrete metric & 0.0850 & 0.1162 & \([0.3182,0.4167]\) & closed form / LP\\
Quadratic distance &  & 0.1333 & \([0.3000,0.4000]\) & closed form / LP\\
Manhattan distance &  & 0.1278 & \([0.3055,0.4050]\) & threshold\\
\hline
\multicolumn{5}{l}{\(n=2,\ k=3,\ r=(1,1,0)\)}\\
Discrete metric & 0.5000 & 0.3889 & \([0.1667,0.5000]\) & closed form / LP\\
Quadratic distance &  & 0.5000 & \([0.1250,0.5000]\) & closed form / LP\\
Manhattan distance &  & 0.4492 & \([0.1464,0.5000]\) & threshold\\
\hline
\multicolumn{5}{l}{\(n=20,\ k=5,\ r=(10,4,3,2,1)\)}\\
Discrete metric & 0.0061 & 0.0714 & \([0.2381,0.3152]\) & closed form / LP\\
Quadratic distance &  & 0.0800 & \([0.2125,0.2875]\) & closed form / LP\\
Manhattan distance &  & 0.0770 & \([0.2182,0.2948]\) & threshold\\
\hline
\end{tabular}
\caption{Fixed-report illustrations. The payment probability is the discrete-metric exact-match probability at \(p=r/n\); it is an implementation cost of that fixed-prize rule, not a property of the other distance rules.}
\label{tab:selected-design}
\end{table}

The fixed-report cases again show why the paper should not claim uniform dominance.
For some reports, discrete-metric scoring gives narrower average coordinate intervals.
For others, quadratic distance gives narrower mean bounds.
The relevant design choice depends on whether the experimenter prioritizes average coordinate precision, worst-coordinate precision, mean inference, implementation robustness, or ease of explanation.

\section{Risk Aversion and Implementation}

Risk aversion matters because a scoring rule can be implemented in different ways.
The distinction is between paying the score directly as money and using the score to determine the probability of a fixed prize.

\paragraph{Direct monetary scores.}
If the subject is paid the score itself, the subject solves
\[
\argmax_{r\in\OmegaN}\E_p[u(S(r,\omega))].
\]
For the frequency-guessing benchmark, this does not change incentives.
The payment is a fixed prize \(B\) if \(r=\omega\) and zero otherwise, so
\[
\E_p[u]
=
u(0)+\Pr_p(\omega=r)\{u(B)-u(0)\}.
\]
Any strictly increasing \(u\) therefore preserves the same optimal reports as maximizing the exact-match probability.

Direct monetary distance scores do not have this robustness in general.
Quadratic and Manhattan distance generate several possible payment levels, so nonlinear utility can change the ranking of reports.
This is the standard risk-aversion problem for monetary scoring rules \cite{ArmantierTreich2012,SchlagEtAl2013}.

\paragraph{Probability implementation.}
The distance rules can instead be implemented as probabilities.
Let \(q(r,\omega)\in[0,1]\) be a normalized score, and pay a fixed prize \(B\) with probability \(q(r,\omega)\), with zero otherwise.
Assume the prize and outside payment are fixed, utility depends only on final money, the randomization is objective and independent, and \(q\) is a valid probability.
Then
\[
\E_p\left[
q(r,\omega)u(B)+(1-q(r,\omega))u(0)
\right]
=
u(0)+(u(B)-u(0))\E_p[q(r,\omega)].
\]
Thus maximizing expected utility is equivalent to maximizing the expected normalized score, independently of the curvature of \(u\).

Because \(\OmegaN\) is finite, each rule considered here is bounded.
A positive affine normalization of a bounded score into \([0,1]\) preserves the expected-score ranking of reports.
Non-affine transformations generally need not preserve the report region.
This is the standard binary-lottery or binarized-scoring-rule logic used to separate incentive rankings from utility curvature \cite{ArmantierTreich2012,SchlagEtAl2013}.
The point is practical rather than novel: direct monetary distance rules are risk-neutral mechanisms, while probability implementations can make them robust to risk aversion under expected utility.
Thus the design comparison above should be read jointly with the implementation choice.
A researcher who values robustness may prefer fixed-prize discrete-metric scoring, while a researcher who wants quadratic-distance belief bounds can use a probability implementation to avoid relying on risk neutrality.

\section{Discussion}

The paper studies frequency-report scoring rules as inferential devices.
The report \(r\) is not the belief vector.
It is evidence about \(p\), and the scoring rule determines how informative that evidence is.
The relevant chain is therefore
\[
S
\rightarrow
R_S(p)
\rightarrow
P_S(r)
\rightarrow
\text{bounds on beliefs and functionals}.
\]

The results suggest three practical recommendations.
First, use discrete-metric frequency scoring when the priority is transparency and fixed-prize robustness.
It has closed-form coordinate bounds and is robust to risk aversion, but exact-match payment probabilities can be low away from concentrated reports.

Second, use quadratic-distance frequency scoring when the priority is transparent belief and mean inference.
It produces a projection rule, a linear belief region, closed-form coordinate bounds, and linear-program bounds for any linear functional.
The design exercise shows that quadratic distance is especially attractive for worst-coordinate uncertainty and for mean-oriented targets in sparse-belief environments.

Third, use Manhattan distance when median-type frequency reports are substantively desired or when it is useful as a comparison rule.
Its binary case is clean, and its multi-category inverse region is explicit, but the coordinate and mean bounds require threshold computation rather than a simple closed form.

The paper should not be read as a general theory of all discrete scoring rules.
Hamming and Chebyshev distance are discussed as additional computable rules, but the primary analytic-bound results are the quadratic-distance and discrete-metric cases.
The intended contribution is narrower: a finite-sample design guide for researchers who want to interpret observed frequency reports as bounds on latent beliefs and target functionals.

Important open questions remain.
The design exercise should be refined before empirical use, especially with alternative outcome vectors, a final choice of paper tables or figures, and robustness checks for tie-breaking conventions.
The cognitive advantage of frequency reports and distance payments should also be tested experimentally rather than assumed.
Those extensions are separate from the main theoretical point: different frequency-report scoring rules induce different inverse belief regions, and those regions determine the informational efficiency of the mechanism.

\appendix

\section{Proofs}

\subsection{Discrete-Metric Frequency Scoring}
\label{app:exact-match-proof}

\begin{proof}[Proof of Proposition 1]
For any report \(r\),
\[
\E_p[a+b\1\{r=\omega\}]
=
a+b\Pr_p(\omega=r),
\]
so expected-score maximization is equivalent to maximizing the multinomial probability mass.

It remains to characterize the mode region.
Compare \(r\) with \(r+e_i-e_j\), which is feasible when \(r_j>0\).
If \(p_j=0\) and \(r_j>0\), then \(r\) cannot be a mode unless all probability is zero, which is impossible.
For \(p_j>0\),
\[
\frac{\Pr_p(r+e_i-e_j)}{\Pr_p(r)}
=
\frac{r_j}{r_i+1}\frac{p_i}{p_j}.
\]
Thus a necessary condition for \(r\) to be a mode is
\[
r_jp_i\le (r_i+1)p_j
\quad
\forall i,j\text{ with }r_j>0.
\]
Conversely, if these inequalities hold, any feasible \(s\) can be reached from \(r\) by a sequence of one-count transfers from coordinates where the current vector exceeds \(s\) to coordinates where it is below \(s\).
At each step the same ratio is at most one, because source counts weakly decrease and destination counts weakly increase along the path.
Therefore no feasible \(s\) has larger multinomial mass than \(r\).

The coordinate upper bound follows by summing
\[
r_jp_i\le (r_i+1)p_j
\]
over \(j\ne i\):
\[
(n-r_i)p_i\le (r_i+1)(1-p_i),
\]
so \(p_i\le (r_i+1)/(n+1)\).
The lower bound follows by summing
\[
r_ip_j\le (r_j+1)p_i
\]
over \(j\ne i\):
\[
r_i(1-p_i)\le (n-r_i+k-1)p_i,
\]
so \(p_i\ge r_i/(n+k-1)\).

The bounds are sharp.
The upper endpoint is attained by
\[
p_i=\frac{r_i+1}{n+1},
\quad
p_j=\frac{r_j}{n+1}\quad(j\ne i).
\]
The lower endpoint, when \(r_i>0\), is attained by
\[
p_i=\frac{r_i}{n+k-1},
\quad
p_j=\frac{r_j+1}{n+k-1}\quad(j\ne i).
\]
When \(r_i=0\), the lower endpoint \(0\) is attained by setting \(p_i=0\) and assigning probabilities proportional to \(r_j\) on the reported support.
When \(r_i=n\), the upper endpoint \(1\) is attained by \(p_i=1\).
\end{proof}

\subsection{Manhattan-Distance Rule}
\label{app:manhattan-proof}

The Manhattan characterization used in Section 3.2 follows from a standard exchange argument.
Linearity of expectation gives
\[
\E_p\left[\sum_{i=1}^k |r_i-\omega_i|\right]
=
\sum_{i=1}^k \E_p|r_i-\omega_i|.
\]
The \(i\)-th term depends only on the binomial marginal \(W_i\sim\mathrm{Bin}(n,p_i)\).
For
\[
\ell_i(t;p_i)=\E|t-W_i|,
\]
the marginal cost of increasing the reported count from \(t\) to \(t+1\) is
\[
\ell_i(t+1;p_i)-\ell_i(t;p_i)
=
2F_i(t;n,p_i)-1.
\]
Moving one reported count from \(j\) to \(i\) changes expected Manhattan loss by
\[
\{2F_i(r_i;n,p_i)-1\}-\{2F_j(r_j-1;n,p_j)-1\}.
\]
Hence no feasible one-count transfer improves the report if and only if
\[
F_i(r_i;n,p_i)\ge F_j(r_j-1;n,p_j)
\quad
\forall i,j\text{ with }r_i<n,\ r_j>0.
\]
Because the marginal increments are weakly increasing in the reported count, this local no-transfer condition is also sufficient for global optimality on the integer simplex.

The threshold representation follows by taking \(c\) between
\[
\max_{j:r_j>0}F_j(r_j-1;n,p_j)
\quad\text{and}\quad
\min_{i:r_i<n}F_i(r_i;n,p_i).
\]
Inverting the resulting binomial-CDF inequalities coordinate by coordinate gives the threshold-indexed box-simplex description used for Manhattan bounds.

\subsection{Quadratic-Distance Rule}
\label{app:squared-proof}

\begin{proof}[Proof of Proposition 2]
Maximizing \(S_2\) is equivalent to minimizing \(\E_p\|r-\omega\|_2^2\).
For fixed \(r\),
\[
\E_p\|r-\omega\|_2^2
=
\|r-\E_p[\omega]\|_2^2
+
\sum_{i=1}^k\mathrm{Var}_p(\omega_i).
\]
The variance term is independent of \(r\), and \(\E_p[\omega]=np\).
Hence optimal reports are exactly the feasible integer projections of \(np\).

The full inverse projection condition is
\[
\|r-np\|_2^2\le \|s-np\|_2^2
\quad
\forall s\in\OmegaN,
\]
or equivalently
\[
2np\cdot(s-r)\le \|s\|_2^2-\|r\|_2^2.
\]
Taking \(s=r+e_i-e_j\), feasible when \(r_j>0\), gives
\[
n(p_i-p_j)\le r_i-r_j+1.
\]
Conversely, suppose all these one-count inequalities hold.
For any \(s\in\OmegaN\), write \(d=s-r\).
Pair each positive unit of \(d_i\) with a negative unit of \(d_j\).
Summing the one-count inequalities over the paired transfers gives
\[
np\cdot d\le r\cdot d+M,
\]
where \(M=\sum_{i:d_i>0}d_i\) is the number of transferred counts.
Because \(d\) has integer components and sums to zero,
\[
M\le \frac12\sum_i d_i^2.
\]
Therefore
\[
2np\cdot(s-r)
\le
2r\cdot(s-r)+\|s-r\|_2^2
=
\|s\|_2^2-\|r\|_2^2,
\]
which is the full projection condition.

It remains to derive the coordinate bounds.
Let \(m=m(r)=|\{j:r_j>0\}|\).
If \(r_i>0\), summing
\[
n(p_u-p_i)\le r_u-r_i+1
\]
over all \(u\ne i\) gives
\[
1-p_i\le (k-1)p_i+\frac{n-kr_i+k-1}{n},
\]
so
\[
p_i\ge \frac{r_i-1}{n}+\frac{1}{nk}.
\]
For the upper bound, sum
\[
n(p_i-p_v)\le r_i-r_v+1
\]
over positive-report coordinates \(v\ne i\), and use nonnegativity of the zero-report coordinates:
\[
1-p_i\ge (m-1)p_i+\frac{n-mr_i-m+1}{n}.
\]
This gives
\[
p_i\le \frac{r_i+1}{n}-\frac{1}{nm}.
\]

If \(r_i=0\), the lower bound is \(0\).
For the upper bound, sum
\[
n(p_i-p_v)\le 1-r_v
\]
over all \(m\) positive-report coordinates \(v\):
\[
1-p_i\ge mp_i+\frac{n-m}{n},
\]
so
\[
p_i\le \frac{m}{n(m+1)}.
\]
The endpoints are attained by making the relevant summed inequalities bind and assigning remaining probability within the categories not used in the binding constraints.
Thus the displayed bounds are sharp.
\end{proof}

\section{Design Exercise Details}
\label{app:design-details}

The design exercise is generated by the script \texttt{scripts/design\_efficiency.py}.
The final grid uses \(n\in\{5,10,20,50\}\), \(k\in\{2,3,5,10\}\), \(\alpha\in\{0.1,0.3,1,3,10\}\), and \(5{,}000\) draws from the symmetric Dirichlet distribution for each cell.
For each draw and rule, the script computes an optimal report and then computes bounds for the inverse region generated by that report.

Discrete-metric coordinate bounds are closed form and mean bounds are linear programs over the transfer region.
Quadratic-distance coordinate bounds are closed form and mean bounds are linear programs over the one-count transfer polytope.
Manhattan coordinate and mean bounds use the threshold representation in Appendix \ref{app:manhattan-proof}; the script inverts binomial CDF inequalities by generalized inverse functions and searches over a threshold grid with tolerance \(10^{-4}\).
The Manhattan entries are therefore threshold-computed.

The generated files are CSV tables under \texttt{outputs/simulation\_design/}.
These outputs are diagnostics for experimental design.
They are not simulations of subject behavior and do not replace the mathematical characterization of \(R_S(p)\) and \(P_S(r)\).


\bibliographystyle{plain}
\bibliography{references}

\end{document}
